% !TEX program = xelatex
% Standalone manuscript. Compile twice with XeLaTeX.
\documentclass[UTF8,a4paper,11pt,fontset=none]{ctexart}
\usepackage{fontspec}
\IfFontExistsTF{Times New Roman}{\setmainfont{Times New Roman}}{\setmainfont{Latin Modern Roman}}
\IfFontExistsTF{Arial}{\setsansfont{Arial}}{\setsansfont{Latin Modern Sans}}
\IfFontExistsTF{Consolas}{\setmonofont{Consolas}}{\setmonofont{Latin Modern Mono}}
\IfFontExistsTF{SimSun}{
  \setCJKmainfont[AutoFakeBold=1.5,ItalicFont=SimSun]{SimSun}
  \setCJKsansfont{Microsoft YaHei}
  \setCJKmonofont{SimSun}
}{
  \setCJKmainfont[BoldFont=FandolSong-Bold.otf,ItalicFont=FandolKai-Regular.otf]{FandolSong-Regular.otf}
  \setCJKsansfont{FandolHei-Regular.otf}
  \setCJKmonofont{FandolFang-Regular.otf}
}
\usepackage[left=25mm,right=25mm,top=23mm,bottom=23mm,headheight=15pt]{geometry}
\usepackage{amsmath,amssymb,amsthm,mathtools,booktabs,array}
\usepackage{xcolor,hyperref,fancyhdr}
\definecolor{linkink}{HTML}{23445E}
\hypersetup{colorlinks=true,linkcolor=linkink,citecolor=linkink,urlcolor=linkink,
  pdftitle={二进制字符串局部模式计数的对数凹性},
  pdfsubject={游程展开、系数平衡与 Cassini 差式}}
\pagestyle{fancy}
\fancyhf{}
\fancyhead[L]{\small 二进制字符串局部模式计数的对数凹性}
\fancyhead[R]{\small 论文稿}
\fancyfoot[C]{\thepage}
\renewcommand{\headrulewidth}{0.3pt}
\setlength{\parindent}{2em}
\setlength{\parskip}{0.20em}
\setlength{\emergencystretch}{2em}
\linespread{1.10}
\allowdisplaybreaks[2]
\numberwithin{equation}{section}
\renewcommand{\proofname}{证明}
\newtheoremstyle{paperplain}{6pt}{6pt}{\normalfont}{}{\bfseries}{.}{0.5em}{}
\theoremstyle{paperplain}
\newtheorem{theorem}{定理}[section]
\newtheorem{lemma}[theorem]{引理}
\newtheorem{proposition}[theorem]{命题}
\newtheorem{corollary}[theorem]{推论}
\theoremstyle{definition}
\newtheorem{definition}[theorem]{定义}
\theoremstyle{remark}
\newtheorem{remark}[theorem]{注}
\DeclareMathOperator{\occ}{occ}
\DeclareMathOperator{\bdes}{bdes}
\DeclareMathOperator{\Impart}{Im}
\newcommand{\lcgap}{D_{n,k}}
\newcommand{\supp}{\operatorname{supp}}
\title{\textbf{二进制字符串局部模式计数的对数凹性}\\[0.3em]
  \large 游程展开、系数平衡与 Cassini 差式}
\author{}
\date{2026年10月14日}

\begin{document}
\maketitle
\begin{abstract}
考虑所有长度为 $n$ 且末位为 $1$ 的二进制字符串，并以连续模式 $10$ 与 $011$
的出现次数之和作为统计量。本文证明，该统计量的计数序列对每个 $n\ge1$ 都对数凹，
且在其正支集的内部严格对数凹。证明从游程分解导出一个非负二项式和，
再用辅助变量记录游程数，将目标差式分成各个固定次数层。
每个高次层的比较分为两步：利用非负线性形式乘积的系数提取证明尾词平衡，
利用 Fibonacci 铺砖多项式的 Cassini 恒等式和单次变号公式证明中央集中。
奇数指标下唯一可能为负的最低项，由紧邻系数的显式下界完全补偿。
文中给出所用多项式事实的直接证明，因而全参数结论不依赖计数多项式本身实根、
额外结构假设或有限计算。
\end{abstract}
\noindent\textbf{关键词：}二进制字符串；局部模式；对数凹；游程；实稳定多项式；Cassini 恒等式。

\section{问题、背景与证明安排}\label{sec:intro}
\subsection{统计量与主定理}
记
\[
\mathcal W_n^{(1)}=\{w=w_1\cdots w_n\in\{0,1\}^n:w_n=1\}.
\]
对一个固定二进制模式 $u$，$\occ_u(w)$ 表示 $u$ 在 $w$ 中作为连续片段出现的次数。
每个起始位置分别计数，允许不同出现彼此重叠。定义
\begin{equation}\label{eq:stat}
s(w)=\occ_{10}(w)+\occ_{011}(w),\qquad
b_{n,k}=\#\{w\in\mathcal W_n^{(1)}:s(w)=k\}.
\end{equation}
约定 $b_{n,k}=0$ 当 $k<0$。例如，$w=1011$ 中两个模式各出现一次，所以 $s(w)=2$。
本文证明以下结论。

\begin{theorem}[主定理]\label{thm:main}
对每个整数 $n\ge1$ 和每个整数 $k$，均有
\begin{equation}\label{eq:main}
b_{n,k}^{\,2}\ge b_{n,k-1}b_{n,k+1}.
\end{equation}
令 $d_n=\lfloor2(n-1)/3\rfloor$。计数序列的正支集恰为
$0\le k\le d_n$，并且当 $1\le k<d_n$ 时，式 \eqref{eq:main} 为严格不等式。
\end{theorem}

\subsection{与模式避免排列的关系}
Elizalde、Rivera 和 Zhuang 的论文 \cite{ERZ} 在 Conjecture 4.6 中提出三阶模式避免排列
的大下降计数序列的对数凹猜想。其 Lemma 3.6、3.9 分别给出
$\mathfrak S_n(123,132)$、$\mathfrak S_n(132,213)$ 到 $\mathcal W_n^{(1)}$ 的双射，
并将大下降数转化为式 \eqref{eq:stat} 的统计量。
因此，定理 \ref{thm:main} 证明了该猜想中这两个等价的分支。
本文不将这一结论扩大为 Conjecture 4.6 的其他分支。

原文已有计数生成函数与递推。本文的论证直接从字符串定义出发，
生成函数只在附录中用于核对；证明的关键在于建立系数不等式，
而不是重新得到已有的递推。

\subsection{证明的组织}
直接证明 $B_n(t)=\sum_kb_{n,k}t^k$ 实根并非本文所需的路线。
我们先引入游程标记多项式 $F_{n,k}(z)$，使 $F_{n,k}(1)=b_{n,k}$，
再研究
\[
D_{n,k}(z)=F_{n,k}(z)^2-F_{n,k-1}(z)F_{n,k+1}(z).
\]
辅助变量 $z$ 的次数记录两条字符串的零游程总数；
它与 $B_n(t)$ 中记录统计量的变量 $t$ 不同。

第 \ref{sec:runs} 节建立游程公式并确定正支集。
第 \ref{sec:tools} 节证明两类需要的基本事实：铺砖多项式的实根与 Cassini 恒等式，
以及非负线性形式乘积的系数提取。
第 \ref{sec:reduction} 节把 $D_{n,k}$ 的固定次数系数写成加权差。
第 \ref{sec:balance}、\ref{sec:concentration} 节分别完成尾词平衡与中央集中，
从而处理所有次数 $R>k$。
第 \ref{sec:exception} 节处理唯一可能出现的最低次负项，
最后在第 \ref{sec:finish} 节合成主定理。

\section{游程展开与支集}\label{sec:runs}
全文对整数上、下标采用约定
\begin{equation}\label{eq:binom-convention}
\binom uv=0\quad\text{除非 }0\le v\le u.
\end{equation}
这也包括上标为负的情况。所有涉及二项式的求和因此都是有限的。

\begin{proposition}[游程公式]\label{prop:runs}
对 $n\ge1$、$k\ge0$，
\begin{equation}\label{eq:runs}
b_{n,k}=\sum_{r\ge0}
\binom{r+1}{k+1-r}\binom{n-r-1}{k}.
\end{equation}
\end{proposition}
\begin{proof}
游程指极大的连续相同符号片段。先考虑首位为 $1$ 的字符串。
若有 $r$ 个零游程，则有 $r+1$ 个一游程；每个零游程前恰有一个 $10$，
所以 $\occ_{10}=r$。除最初的一游程外，设有 $q$ 个一游程的长度至少为二。
它们各贡献一次 $011$，故 $k=r+q$。

选出这 $q$ 个一游程有 $\binom rq=\binom r{k-r}$ 种方法。
最小总长度为 $1+2r+q=1+r+k$。
额外长度只能加到最初的一游程、$r$ 个零游程或 $q$ 个长一游程，
共 $k+1$ 个位置；其他一游程必须保持长度一。
当额外长度非负时，非负整数分配的数量为
\[
\binom{n-r-k-1+k}{k}=\binom{n-r-1}{k}.
\]
当额外长度为负时，式 \eqref{eq:binom-convention} 使数量为零。

再考虑首位为 $0$ 的字符串。此时 $r\ge1$ 个零游程对应 $r$ 个一游程，
$\occ_{10}=r-1$。若 $q$ 个一游程的长度至少为二，则 $k=r-1+q$。
选择它们有 $\binom r{k+1-r}$ 种方法。
最小总长度同样为 $k+r+1$，额外长度有 $r+q=k+1$ 个分配位置，
因此所得数量为
\[
\binom r{k+1-r}\binom{n-r-1}{k}.
\]
将两种情形相加，使用 Pascal 恒等式
$\binom r{k-r}+\binom r{k+1-r}=\binom{r+1}{k+1-r}$，即得 \eqref{eq:runs}。
$r=0$ 的项只来自首位为 $1$ 的情形；$n=1$ 时该项为一，
其余项为零，所以没有遗漏最短字符串。
\end{proof}

\begin{corollary}[正支集与非平凡范围]\label{cor:support}
对 $k\ge0$，
\begin{equation}\label{eq:support}
b_{n,k}>0\quad\Longleftrightarrow\quad
n\ge k+\lceil k/2\rceil+1.
\end{equation}
因而正支集为 $0\le k\le d_n$。若 $k\ge1$ 且 $b_{n,k+1}>0$，则
\begin{equation}\label{eq:range}
n\ge k+2+\lceil(k+1)/2\rceil\ge k+3.
\end{equation}
\end{corollary}
\begin{proof}
式 \eqref{eq:runs} 的第一因子非零恰要求
$\lceil k/2\rceil\le r\le k+1$，第二因子非零要求 $n\ge r+k+1$。
必要性随即成立；反过来，取 $r=\lceil k/2\rceil$ 就得到一个正项。
条件 $k+\lceil k/2\rceil\le n-1$ 等价于 $k\le\lfloor2(n-1)/3\rfloor$。
以 $k+1$ 代替 $k$，得到最后的结论。
\end{proof}

式 \eqref{eq:range} 有一个重要用途：
若右邻项 $b_{n,k+1}$ 为零，对数凹不等式已经成立。
所以所有需要实质证明的情形都落在 $k\ge1$、$n\ge k+3$ 内；
后面的倒易多项式因而不会产生负幂。

\section{辅助多项式与自足的系数工具}\label{sec:tools}
\subsection{铺砖多项式及其差式}
对 $j\ge0$ 定义
\begin{equation}\label{eq:Q}
q_{j,r}=\binom{r+1}{j+1-r},\qquad
Q_j(z)=\sum_{r\ge0}q_{j,r}z^r.
\end{equation}
第一因子还具有铺砖意义：用长度一、二的砖铺长度 $j+2$ 的板，
若砖数为 $r+1$，则长度二的砖数为 $j+1-r$，
相应铺法数正是 $q_{j,r}$。

\begin{lemma}[实根与 Cassini 恒等式]\label{lem:Q}
$Q_j$ 的次数为 $j+1$，最小次数为 $\lceil j/2\rceil$，并且
\begin{equation}\label{eq:Q-factor}
Q_j(z)=z^{\lceil j/2\rceil}
\prod_{a=1}^{\lfloor(j+2)/2\rfloor}
\left(z+4\cos^2\frac{a\pi}{j+3}\right).
\end{equation}
特别地，其全部根为非正实数。对 $k\ge1$，
\begin{equation}\label{eq:cassini}
Q_k(z)^2-Q_{k-1}(z)Q_{k+1}(z)=(-z)^k.
\end{equation}
\end{lemma}
\begin{proof}
令
\[
H_L(x)=\sum_{a=0}^{\lfloor L/2\rfloor}\binom{L-a}{a}x^a.
\]
按第一块砖的长度分解，$H_0=H_1=1$，
$H_L=H_{L-1}+xH_{L-2}$。
对 $0<\theta<\pi/2$，代入 $x=-1/(4\cos^2\theta)$，
用初值和递推直接核验
\[
H_L\left(-\frac1{4\cos^2\theta}\right)
=\frac{\sin((L+1)\theta)}{(2\cos\theta)^L\sin\theta}.
\]
取 $\theta=a\pi/(L+1)$，$1\le a\le\lfloor L/2\rfloor$，
得到 $\lfloor L/2\rfloor$ 个不同的负根。
这正是 $H_L$ 的次数，因此已经列出全部根。
其常数项为一，于是
\[
H_L(x)=\prod_{a=1}^{\lfloor L/2\rfloor}
\left(1+4x\cos^2\frac{a\pi}{L+1}\right).
\]
将求和指标反向可得 $Q_j(z)=z^{j+1}H_{j+2}(1/z)$，从而得
\eqref{eq:Q-factor}，包括零根的重数。

另一方面，$Q_0=1+z$、$Q_1=2z+z^2$，
且 $Q_j=z(Q_{j-1}+Q_{j-2})$ 对 $j\ge2$ 成立。
记 $\Delta_k=Q_k^2-Q_{k-1}Q_{k+1}$，则 $\Delta_1=-z$，
而递推给出
\[
\Delta_k=z(Q_kQ_{k-2}-Q_{k-1}^2)=-z\Delta_{k-1}\qquad(k\ge2).
\]
归纳即得 \eqref{eq:cassini}。
\end{proof}

\subsection{系数提取为何保留需要的性质}
接下来的工具只作用于非负线性形式的乘积，
不需要假设 $B_n(t)$ 实根，也不需要调用一般的稳定化定理。

\begin{definition}
实系数二元多项式 $P(X,Y)$ 称为实稳定，
若 $\Impart X>0$、$\Impart Y>0$ 时 $P(X,Y)\ne0$。
零多项式单独处理，不包含在此定义中。
\end{definition}

\begin{lemma}[线性形式乘积的系数提取]\label{lem:extract}
设
\[
G(X,Y,Z)=c\prod_{i=1}^M(\alpha_iX+\beta_iY+\gamma_iZ),
\]
其中 $c>0$，$\alpha_i,\beta_i,\gamma_i\ge0$ 且 $\alpha_i+\beta_i>0$。
则对每个非负整数 $a$，$[Z^a]G$ 为零或为实稳定多项式。
\end{lemma}
\begin{proof}
固定上半平面的 $X,Y$。记 $d=\#\{i:\gamma_i>0\}$。
作为 $Z$ 的多项式，$G$ 的最高系数为
\[
c\left(\prod_{\gamma_i>0}\gamma_i\right)
\left(\prod_{\gamma_i=0}(\alpha_iX+\beta_iY)\right)\ne0.
\]
所以它的次数恒为 $d$，与所选 $X,Y$ 无关。
当 $d>0$ 时，其根
$-(\alpha_iX+\beta_iY)/\gamma_i$ 都严格位于下半平面。

若非常数多项式 $p$ 的全部根 $\lambda_i$ 在下半平面，
则在 $\Impart Z\ge0$ 上
\[
\frac{p'(Z)}{p(Z)}=\sum_i\frac1{Z-\lambda_i}
\]
的虚部严格为负。因此 $p'$ 在该闭半平面也没有根。
反复求导，直到非零常数阶段，这一性质一直保持。
当 $0\le a\le d$ 时，$p^{(a)}(0)\ne0$；
当 $a>d$ 时导数恒为零。
由 $[Z^a]G=p^{(a)}(0)/a!$，所需结论成立。
\end{proof}

\begin{lemma}[二元齐次因式与系数对数凹]\label{lem:homogeneous}
设 $P(X,Y)$ 非零、齐次、实稳定，且所有系数非负。
则存在 $c>0$、非负整数 $\alpha,\beta$ 及正数 $a_i,b_i$，使
\begin{equation}\label{eq:binary-factor}
P(X,Y)=cX^\alpha Y^\beta\prod_i(a_iX+b_iY).
\end{equation}
此外，若 $P=\sum_{j=0}^d p_jX^jY^{d-j}$，则 $(p_j)$ 对数凹且没有内部零。
\end{lemma}
\begin{proof}
若 $P(t,1)$ 有上半平面根 $t$，
选择幅角 $\varepsilon>0$ 充分小的 $Y$，使 $X=tY$ 仍在上半平面。
齐次性给出 $P(X,Y)=Y^dP(t,1)=0$，与稳定性矛盾。
实系数也排除下半平面根，非负系数排除正实根。
提出零根及纯 $Y$ 因子，即得 \eqref{eq:binary-factor}。

下面直接证明其系数对数凹性。若非负序列 $(u_j)$ 对数凹且无内部零，
令 $v_j=a u_j+b u_{j-1}$，其中 $a,b\ge0$，则
\begin{align*}
v_j^2-v_{j-1}v_{j+1}
={}&a^2(u_j^2-u_{j-1}u_{j+1})
+b^2(u_{j-1}^2-u_{j-2}u_j)\\
&+ab(u_ju_{j-1}-u_{j-2}u_{j+1})\ge0.
\end{align*}
前两项非负。若最后一项的右侧乘积非零，四项均处于连续的正支集内，
由相邻比递减得到 $u_ju_{j-1}\ge u_{j-2}u_{j+1}$；
否则该比较直接成立。
乘一个非负线性因子也保持支集连续。
从单项式起逐个乘因子，就得到结论。
\end{proof}

\section{把目标差式分解为固定次数层}\label{sec:reduction}
依据游程公式，定义
\begin{equation}\label{eq:F}
F_{n,j}(z)=\sum_{r\ge0}q_{j,r}\binom{n-r-1}{j}z^r,\qquad
F_{n,j}(1)=b_{n,j}.
\end{equation}
对固定 $k$，令
\begin{equation}\label{eq:D}
D_{n,k}(z)=F_{n,k}(z)^2-F_{n,k-1}(z)F_{n,k+1}(z).
\end{equation}
我们最终只需证明 $D_{n,k}(1)\ge0$。
不过，将 $z$ 的次数逐层固定，会将两条不同统计量的比较拆成两个可处理的比较。

固定 $k\ge1$、$n\ge k+3$ 和整数 $R>k$。
以下求和均为 $0\le r\le R$，并始终记 $s=R-r$。
定义
\begin{align}
a_r&=q_{k,r}q_{k,s},&
c_r&=q_{k-1,r}q_{k+1,s},\label{eq:ac}\\
W_r&=\binom{n-r-1}{k}\binom{n-s-1}{k},&
U_r&=\binom{n-r-1}{k-1}\binom{n-s-1}{k+1}.\label{eq:WU}
\end{align}
于是
\begin{equation}\label{eq:layer}
[z^R]D_{n,k}=\sum_r a_rW_r-\sum_r c_rU_r.
\end{equation}
其中 $a_r,c_r$ 来自铺砖部分，$W_r,U_r$ 来自额外长度分配，
亦可理解为在两个尾词中选定相应数量的位置。

证明的核心是以下两步链条：
\begin{equation}\label{eq:chain}
\sum_r c_rU_r
\ \le\ \sum_r c_rW_r
\ \le\ \sum_r a_rW_r.
\end{equation}
第一步保持铺砖系数 $c_r$，将尾词的选位数 $k-1,k+1$ 平衡为 $k,k$。
第二步保持平衡后的尾词权重 $W_r$，比较两个铺砖分布。
这两步分别在后面两节证明，彼此不以原对数凹命题为前提。

\section{第一步比较：尾词平衡}\label{sec:balance}
\begin{proposition}\label{prop:balance}
在式 \eqref{eq:ac}、\eqref{eq:WU} 的记号与范围内，
\begin{equation}\label{eq:balance}
\sum_r c_rU_r\le\sum_r c_rW_r.
\end{equation}
\end{proposition}
\begin{proof}
若 $R>2k+2$，由 $\deg Q_{k-1}=k$、$\deg Q_{k+1}=k+2$，
所有 $c_r$ 都为零。以下设 $R\le2k+2$。

\textbf{先构造具有实稳定性的固定次数核。}
令
\[
C_R(X,Y)=\sum_r c_rX^rY^s.
\]
把 $Q_j$ 齐次化至其次数 $j+1$：
\[
\widetilde Q_j(X,Z)=\sum_rq_{j,r}X^rZ^{j+1-r}.
\]
由式 \eqref{eq:Q-factor}，这是单项式与正系数线性形式的乘积。
并且
\[
C_R(X,Y)=[Z^{2k+2-R}]
\bigl(\widetilde Q_{k-1}(X,Z)\widetilde Q_{k+1}(Y,Z)\bigr).
\]
每个因子都符合引理 \ref{lem:extract} 的条件，
所以 $C_R$ 为零或实稳定。若为零，命题平凡；以下设其非零。

\textbf{再把尾词的二项式权重装入同一个多项式。}
由引理 \ref{lem:homogeneous}，写
$C_R=cX^\alpha Y^\beta\prod_i(a_iX+b_iY)$。
其中纯 $X$、纯 $Y$ 因子已经提出，故 $a_i,b_i>0$。
因为 $\deg_X C_R\le k$、$\deg_Y C_R\le k+2$，
以及 $n-1\ge k+2$，倒易式
\begin{align}
T_R(X,Y)
&=X^{n-1}Y^{n-1}C_R(X^{-1},Y^{-1})\notag\\
&=\sum_r c_rX^{n-r-1}Y^{n-s-1}\label{eq:T}\\
&=cX^{n-1-\deg_X C_R}Y^{n-1-\deg_Y C_R}
\prod_i(a_iY+b_iX)\notag
\end{align}
的所有幂次均非负。
因此 $Z_R(x,y)=T_R(1+x,1+y)$ 是非负仿射线性形式的乘积，
每个因子的常数项均为正。

其总次数为 $d_0=2n-2-R$。
令总次数 $2k$ 的齐次部分为
\[
K_R(x,y)=\sum_{j=0}^{2k}h_jx^jy^{2k-j}.
\]
若 $d_0<2k$，则全部 $h_j$ 为零，命题已成立。
否则，把 $Z_R$ 每个仿射因子的常数项乘共同变量 $w$，
再提取 $w^{d_0-2k}$ 的系数，所得正是 $K_R$。
再次使用引理 \ref{lem:extract}，$K_R$ 为零或实稳定。
排除平凡的零情形，引理 \ref{lem:homogeneous} 给出
\begin{equation}\label{eq:h-lc}
h_k^2\ge h_{k-1}h_{k+1}.
\end{equation}

\textbf{最后利用铺砖分布的方向偏置，确定中央系数的大小。}
当 $r\le s$ 时，
\begin{equation}\label{eq:c-bias}
c_r\ge c_s.
\end{equation}
事实上，在共同正支集上，
\[
\frac{q_{k+1,r}}{q_{k-1,r}}
=\frac{(2r+1-k)(2r-k)}{(k+1-r)(k+2-r)}
\]
随 $r$ 递增。
两列的正支集分别为
$[\lceil(k-1)/2\rceil,k]$ 与 $[\lceil(k+1)/2\rceil,k+2]$。
若交叉积 $q_{k-1,s}q_{k+1,r}$ 正且 $r\le s$，
则 $r,s$ 必在共同正支集，所以上述比值证明
$q_{k-1,r}q_{k+1,s}\ge q_{k-1,s}q_{k+1,r}$；
若右侧为零，该比较直接成立。

记 $A=n-r-1$、$B=n-s-1$。
当 $r\le s$ 时 $A\ge B$，并有
\begin{equation}\label{eq:binom-order}
\binom A{k+1}\binom B{k-1}
\ge\binom A{k-1}\binom B{k+1}.
\end{equation}
这由比值
$(A-k+1)(A-k)/(k(k+1))$ 在分母非零的范围 $A\ge k-1$ 内递增得到；
其他情况按零支集直接检验。
将 $r,s$ 配对，式 \eqref{eq:c-bias}、\eqref{eq:binom-order} 给出
\begin{align*}
h_{k+1}-h_{k-1}
={}&\sum_{r<s}(c_r-c_s)
\left(\binom A{k+1}\binom B{k-1}
-\binom A{k-1}\binom B{k+1}\right)\ge0.
\end{align*}
结合 \eqref{eq:h-lc}，非负性使 $h_k\ge h_{k-1}$。
由系数的定义，
\[
h_k=\sum_r c_rW_r,\qquad h_{k-1}=\sum_r c_rU_r.
\]
因此 \eqref{eq:balance} 成立。
\end{proof}

这一步并非对每个 $r$ 逐项比较 $U_r$ 与 $W_r$；
所需不等式来自固定次数核的整体实稳定性和反射偏置。
第二步将使用不同的机制，只比较已经对称的权重 $W_r$。

\section{第二步比较：中央集中}\label{sec:concentration}
\begin{proposition}\label{prop:concentration}
在相同记号与范围内，
\begin{equation}\label{eq:concentration}
\sum_r c_rW_r\le\sum_r a_rW_r.
\end{equation}
\end{proposition}
\begin{proof}
$R>2k+2$ 时所有项为零。以下设 $k<R\le2k+2$。
令 $c_r^*=(c_r+c_s)/2$。因为 $a_r=a_s$、$W_r=W_s$，
可以把待比较的差写为
\[
\sum_r(a_r-c_r^*)W_r.
\]
我们先确定这个对称差列的符号，再用权重的中央单调性求和。

设
\begin{align*}
d&=r-s,&p&=R-k,&L&=2k+2-R,\\
A&=p+1,&B&=L+2,&P&=AB,&S&=A+B=k+5,\\
H&=2A+B-2=R+4,&J&=2S-3=2k+7.
\end{align*}
其中 $p\ge1$、$L\ge0$。
$a_r>0$ 的支集是中央对称整数区间，在其中
$|d|\le\min(p,L)$，故 $A\pm d>0$、$B\pm d>0$。
用二项式相邻比计算得
\begin{align}
\frac{c_r}{a_r}
&=\frac{(L-d)(p-d)}{(p+1+d)(L+2+d)}\notag\\
&=1-\frac{H+Jd}{P+Sd+d^2}.\label{eq:ratio}
\end{align}
例如，其中第一因子的比为
$(k+1-r)/(2r+1-k)$，第二因子的比为
$(2s-k)/(k+2-s)$；代入 $r=(R+d)/2$ 即得到首行。
在支集端点某个分子为零时，同一公式仍成立。

把 \eqref{eq:ratio} 与其 $d\mapsto-d$ 的版本平均，得到
\begin{equation}\label{eq:single-cross}
\frac{c_r^*}{a_r}-1
=\frac{-HP+(JS-H)d^2}{(P+d^2)^2-S^2d^2}.
\end{equation}
其分母为
$(A+d)(A-d)(B+d)(B-d)>0$，而
\[
JS-H=2k^2+17k+31-R>0.
\]
因此 $a_r-c_r^*$ 的符号随 $|d|$ 增大至多从正变负一次。
在 $a_r=0$ 的外侧，其差本来就非正。
这证明差列在中央非负、外侧非正，不会再次变为正。

Cassini 恒等式 \eqref{eq:cassini} 给出零总质量：
\begin{equation}\label{eq:zero-mass}
\sum_r(a_r-c_r^*)
=[z^R](Q_k^2-Q_{k-1}Q_{k+1})=0,
\end{equation}
因为 $R>k$。
另一方面，$W_r$ 对称且随 $|r-s|$ 不增。
为明确这一点，若 $r<s$ 且 $W_r>0$，
记 $A_0=n-r-1\ge B_0+1$、$B_0=n-s-1\ge k$，则向中央移动一步后
\[
\frac{W_{r+1}}{W_r}
=\frac{A_0-k}{A_0}\frac{B_0+1}{B_0+1-k}\ge1,
\]
最后一个不等式等价于 $k(A_0-B_0-1)\ge0$。
若 $W_r=0$，向中央移动不会违反非负单调性。

若差列恒为零，则结论已得。
否则，由 \eqref{eq:zero-mass}，它有正项和负项。
单次变号性质使所有正项比所有负项更靠近中央，
所以可以选取 $w_0$，使正项处 $W_r\ge w_0$、负项处 $W_r\le w_0$。
因此
\begin{align*}
\sum_r(a_r-c_r^*)W_r
&=\sum_r(a_r-c_r^*)(W_r-w_0)
+w_0\sum_r(a_r-c_r^*)\\
&\ge0.
\end{align*}
最后使用 $\sum c_r^*W_r=\sum c_rW_r$，得到 \eqref{eq:concentration}。
\end{proof}

\begin{corollary}[所有高次层非负]\label{cor:high}
对 $k\ge1$、$n\ge k+3$ 和 $R>k$，
\begin{equation}\label{eq:high-positive}
[z^R]D_{n,k}(z)\ge0.
\end{equation}
\end{corollary}
\begin{proof}
命题 \ref{prop:balance}、\ref{prop:concentration} 给出链条 \eqref{eq:chain}，
代回 \eqref{eq:layer} 即可。
\end{proof}

至此，除最低次数外的全部层已经处理完毕。
不能直接宣称 $D_{n,k}$ 的全部系数非负：
Cassini 恒等式在奇数 $k$ 时确实留下一个负项。
下面明确计算这一异常，并证明下一项足以补偿它。

\section{最低次数与奇数指标的补偿}\label{sec:exception}
\subsection{最低两层的精确系数}
引理 \ref{lem:Q} 表明 $F_{n,j}$ 不含次数小于 $\lceil j/2\rceil$ 的项。
因此，$R<k$ 时 $[z^R]D_{n,k}=0$。
若 $k=2m$ 为偶数，$F_{n,k-1}F_{n,k+1}$ 的最小可能次数为 $k+1$，
而
\begin{equation}\label{eq:even-low}
[z^k]D_{n,k}=\binom{n-m-1}{k}^2\ge0.
\end{equation}
由推论 \ref{cor:high}，偶数情形已经完全解决。

设 $k=2m+1$ 为奇数，记 $N=n-m-2$。
最低层只来自乘积 $F_{n,k-1}F_{n,k+1}$，故
\begin{equation}\label{eq:odd-low}
[z^k]D_{n,k}=-A_0,\qquad
A_0=\binom{N+1}{k-1}\binom N{k+1}.
\end{equation}
为计算紧邻层，所需的铺砖系数为
\begin{align*}
q_{k,m+1}&=m+2,\\
q_{k-1,m}&=q_{k+1,m+1}=1,\\
q_{k-1,m+1}&=\binom{m+2}{2},&
q_{k+1,m+2}&=\binom{m+3}{2}.
\end{align*}
于是
\begin{align}
E:=[z^{k+1}]D_{n,k}
={}&(m+2)^2\binom Nk^2\notag\\
&-\binom{m+3}{2}\binom{N+1}{k-1}\binom{N-1}{k+1}\notag\\
&-\binom{m+2}{2}\binom N{k-1}\binom N{k+1}.\label{eq:E}
\end{align}

\subsection{完全补偿的显式界}
\begin{proposition}\label{prop:compensation}
若式 \eqref{eq:odd-low} 中 $A_0>0$，则 $E>A_0$。
更精确地，令 $C=\binom Nk$，则
\begin{equation}\label{eq:gap-bound}
E-A_0\ge
\frac{2m^2+8m+17}{6(k+1)}C^2>0.
\end{equation}
\end{proposition}
\begin{proof}
$A_0>0$ 要求 $N\ge k+1$。
令 $h=N-k\ge1$，则 $C>0$。
式 \eqref{eq:E} 中两种负乘积分别满足
\begin{align}
\frac{\binom{N+1}{k-1}\binom{N-1}{k+1}}{C^2}
&=\frac{k}{k+1}
\frac{(N+1)h(h-1)}{N(h+1)(h+2)}
\le\frac{k}{k+1},\label{eq:negative-one}\\
\frac{\binom N{k-1}\binom N{k+1}}{C^2}
&=\frac{k}{k+1}\frac h{h+1}
\le\frac{k}{k+1}.\label{eq:negative-two}
\end{align}
其中 \eqref{eq:negative-one} 的不等式由
\[
N(h+1)(h+2)-(N+1)h(h-1)
=3h^2+(4k+3)h+2k>0
\]
得到。利用
$\binom{m+3}{2}+\binom{m+2}{2}=(m+2)^2$，有
\begin{equation}\label{eq:E-bound}
E\ge\frac{(m+2)^2}{k+1}C^2.
\end{equation}

另一方面，令整数 $\ell=h+1\ge2$。直接相除得
\begin{align}
\frac{A_0}{C^2}
&=\frac{k}{k+1}
\frac{(\ell+k)(\ell-1)}{\ell(\ell+1)}\notag\\
&=\frac{k}{k+1}
\left(\frac{\ell-1}{\ell+1}
+k\frac{\ell-1}{\ell(\ell+1)}\right)
\le\frac{k}{k+1}\left(1+\frac k6\right).\label{eq:A-bound}
\end{align}
这里 $\frac{\ell-1}{\ell+1}\le1$；
又因 $\ell$ 是不小于二的整数，
$(\ell-2)(\ell-3)\ge0$，所以
$\frac{\ell-1}{\ell(\ell+1)}\le\frac16$。
这一界使用了整数性，并非对所有实数 $\ell\ge2$ 作同样断言。

最后，$k=2m+1$ 给出
\[
(m+2)^2-k-\frac{k^2}{6}
=\frac{2m^2+8m+17}{6}>0.
\]
从 \eqref{eq:E-bound} 减去 \eqref{eq:A-bound} 即得 \eqref{eq:gap-bound}。
\end{proof}

若 $A_0=0$，最低层并非负项，高次层非负性已经足够。
若 $A_0>0$，命题 \ref{prop:compensation} 表明最低两层的系数之和严格为正。
因此无论哪种情形，奇数指标都不留下未处理的异常。

\section{主定理的合成与排列推论}\label{sec:finish}
\begin{proof}[定理 \ref{thm:main} 的证明]
正支集已由推论 \ref{cor:support} 确定。
若 $k\le0$ 或 $b_{n,k+1}=0$，不等式直接成立。
其余情形由 \eqref{eq:range} 落在前述各命题适用的范围内。

偶数 $k$ 时，次数小于 $k$ 的系数为零，
最低层由 \eqref{eq:even-low} 非负，其余层由 \eqref{eq:high-positive} 非负。
因此 $D_{n,k}(1)\ge0$。
奇数 $k$ 时，次数小于 $k$ 的系数仍为零，
最低层与下一层之和为 $E-A_0\ge0$，
剩余层亦非负，所以仍有 $D_{n,k}(1)\ge0$。
结合 \eqref{eq:F}、\eqref{eq:D}，
\[
D_{n,k}(1)=b_{n,k}^2-b_{n,k-1}b_{n,k+1},
\]
对数凹性得证。

若 $1\le k<d_n$，则右邻项为正。
偶数 $k=2m$ 时，\eqref{eq:support} 对 $k+1$ 给出 $n\ge3m+3$，
所以 \eqref{eq:even-low} 严格为正。
奇数 $k=2m+1$ 时，同一支集条件给出 $n\ge3m+4$，
故 $N\ge k+1$、$A_0>0$，命题 \ref{prop:compensation} 给出严格的正余量。
因此正支集内部的全部不等式都是严格的。
\end{proof}

\begin{corollary}[两个等价的排列分支]\label{cor:permutations}
对 $n\ge1$，集合 $\mathfrak S_n(123,132)$ 与 $\mathfrak S_n(132,213)$
上的大下降数计数序列均对数凹。
\end{corollary}
\begin{proof}
大下降指满足 $\pi_i>\pi_{i+1}+1$ 的位置。
由 \cite[Lemmas 3.6 and 3.9]{ERZ} 的统计保持双射，
这两个计数序列均逐项等于 $(b_{n,k})$。
定理 \ref{thm:main} 因而直接适用。
\end{proof}

整个证明使用的逻辑链为：
游程分解建立加权系数公式，
铺砖实根与两次系数提取建立尾词平衡，
Cassini 零总质量与单次变号建立中央集中，
两步比较处理所有高次层，
显式补偿处理唯一最低异常。
各环节的结论均在应用前证明，不存在向另一个未证对数凹命题的归约。

\appendix
\section{生成函数、递推与小规模核对}\label{sec:appendix}
本附录仅提供与已有计数结构的衔接，不参与主定理的逻辑依赖。
令 $B_n(t)=\sum_kb_{n,k}t^k$，
$G(x,t)=\sum_{n\ge1}B_n(t)x^n$。
正长度零游程的生成函数为 $Z=x/(1-x)$。
一个紧随零游程的一游程，长度一时贡献 $x$，长度至少为二时贡献
$tx^2/(1-x)$，故记
$A=x+tx^2/(1-x)$。
首位为一时最初的一游程贡献 $x/(1-x)$；
首位为零时最初的一对游程贡献 $ZA$。
每个后续的零、一游程对额外贡献一次 $10$，故其权重为 $tZA$。
因此
\begin{align}
G(x,t)
&=\frac{x/(1-x)+ZA}{1-tZA}\notag\\
&=\frac{x\{1+(t-1)x^2\}}
{1-2x+(1-t)x^2+t(1-t)x^3}.\label{eq:GF}
\end{align}
比较 $x^n$ 系数，得到
\[
B_n(t)=2B_{n-1}(t)+(t-1)B_{n-2}(t)
+t(t-1)B_{n-3}(t)\qquad(n\ge4),
\]
初值为 $B_1=1$、$B_2=2$、$B_3=2+2t$。
例如
\[
B_4=2+5t+t^2,\qquad B_5=2+8t+6t^2.
\]
这些式子与游程公式一致，也与 \cite{ERZ} 中的已有生成函数相容。

正文对所有 $n,k,R$ 的证明由所列恒等式、不等式及支集边界独立组成；
上述初值和小规模例子只用于核对计数定义，不承担全参数论证。
配套 LaTeX 源文件为独立文档，可用 XeLaTeX 两遍编译，
无需额外图片、数据文件或外部数学计算系统。

\begin{thebibliography}{9}
\bibitem{ERZ}
Sergi Elizalde, Johnny Rivera Jr., and Yan Zhuang.
\newblock Counting Pattern-Avoiding Permutations by Big Descents.
\newblock \emph{Annals of Combinatorics}, 2026.
\newblock DOI: \href{https://doi.org/10.1007/s00026-026-00805-1}{10.1007/s00026-026-00805-1}.
\end{thebibliography}
\end{document}
